% =====================================================================
%  setup/layout.tex -- settings specific to the A2 poster (AE3)
%  The shared code is in setup/fyp-common.tex.
%
%  Brand: Southampton Solent University brand guidelines (v2.0, 2024):
%  Solent Red and white as lead colours, Navy for text, logo top left,
%  recommended logo size for A2 about 60 mm (A3: 45 mm, minimum 40 mm).
% =====================================================================
\makeatletter
\def\fyp@unit{poster}
\def\fyp@doctype{Poster}
\providecommand\posterorientation{landscape}
\edef\fyp@orient{\posterorientation}\edef\fyp@portraitstr{portrait}
\newif\iffyp@portrait
\ifx\fyp@orient\fyp@portraitstr \fyp@portraittrue \fi
\makeatother

% =====================================================================
%  setup/fyp-common.tex -- shared code of the SEF639 LaTeX templates
%
%  This file is IDENTICAL in all templates (Progress Report, Final
%  Report, Poster). Its master copy lives in shared/setup/ in the template
%  repository; tools/sync-shared.sh copies it into each template.
%  Do not edit the copies -- edit the master and run the script.
%
%  Students: you should not need to edit this file. Your own details
%  belong in main.tex, your own packages at the end of its preamble.
%
%  Loaded by setup/layout.tex, which sets before loading:
%    \fyp@unit     chapter | section | poster  (unit of the main body)
%    \fyp@doctype  e.g. Final Report
% =====================================================================
\makeatletter
\providecommand\fyp@unit{chapter}
\providecommand\fyp@doctype{Report}
\def\fyp@chapterstr{chapter}
\def\fyp@posterstr{poster}
\newif\iffyp@poster
\newif\iffyp@maxmode     % word count: maximum instead of target
\newif\iffyp@guidemode   % word count: recommendation only
\ifx\fyp@unit\fyp@chapterstr
  \def\fyp@unitname{Chapter}\let\fyp@addunnumbered\addchap
\else
  \def\fyp@unitname{Section}\let\fyp@addunnumbered\addsec
\fi
\ifx\fyp@unit\fyp@posterstr \fyp@postertrue \def\fyp@unitname{Panel}\fi

% ---------------------------------------------------------------------
% Draft or final mode (set \reportmode in main.tex)
% ---------------------------------------------------------------------
\newif\iffinalmode
\edef\fyp@mode{\reportmode}\edef\fyp@finalstr{final}
\ifx\fyp@mode\fyp@finalstr \finalmodetrue \fi

% ---------------------------------------------------------------------
% Language and fonts (LuaLaTeX)
% ---------------------------------------------------------------------
% \fontstyle (set at the top of main.tex):
%   brand   -- Carlito, a free font with the same metrics as Calibri, the
%              system font of the Southampton Solent University brand
%              guidelines (the brand fonts Bliss and Bilo are licensed)
%   classic -- Latin Modern, the standard LaTeX look (serif)
\providecommand\fontstyle{brand}
\usepackage{fontspec}
\edef\fyp@font{\fontstyle}\edef\fyp@brandstr{brand}
\ifx\fyp@font\fyp@brandstr
  \setmainfont{Carlito}[Extension=.ttf, UprightFont=*-Regular,
    BoldFont=*-Bold, ItalicFont=*-Italic, BoldItalicFont=*-BoldItalic]
  \setsansfont{Carlito}[Extension=.ttf, UprightFont=*-Regular,
    BoldFont=*-Bold, ItalicFont=*-Italic, BoldItalicFont=*-BoldItalic]
\fi
\usepackage{polyglossia}
\setdefaultlanguage[variant=british]{english}
\usepackage[autostyle]{csquotes}
\usepackage{microtype}

% ---------------------------------------------------------------------
% Page layout
% ---------------------------------------------------------------------
\iffyp@poster\else     % the poster sets its own page size
\usepackage[a4paper, top=25mm, bottom=25mm, left=25mm, right=25mm,
  includehead, includefoot, marginparwidth=20mm, marginparsep=2mm]{geometry}
\usepackage{setspace}
\setstretch{1.15}                % slightly more open than single spacing
\fi

% ---------------------------------------------------------------------
% Colours: Southampton Solent University brand guidelines (v2.0, 2024)
% Lead colours are Solent Red and white, with Navy for text and
% headings. Tints: only 80%, 60% and 40% of a colour (e.g. SkyBlue!40);
% any tint of black, but not below 10% (black!10) in print.
% ---------------------------------------------------------------------
\usepackage[table]{xcolor}
\definecolor{SolentRed}{HTML}{CD1423}
\definecolor{Navy}{HTML}{1D2136}
\definecolor{SkyBlue}{HTML}{01C3DE}
\definecolor{SeaBlue}{HTML}{0F435F}
\definecolor{Teal}{HTML}{1F988A}
\definecolor{PetrolBlue}{HTML}{083237}
\definecolor{Orange}{HTML}{FF4D00}
\definecolor{Yellow}{HTML}{FFC11B}
\definecolor{Green}{HTML}{C6C745}
\definecolor{Purple}{HTML}{9243F4}
\definecolor{Pink}{HTML}{D82F89}
\definecolor{Burgundy}{HTML}{531623}
% risk ratings (risk register): 40% tints of brand colours
\colorlet{RiskLow}{Green!40}
\colorlet{RiskMedium}{Yellow!40}
\colorlet{RiskHigh}{Orange!40}
\colorlet{RiskVeryHigh}{SolentRed!40}

% ---------------------------------------------------------------------
% Figures, tables, lists, units, maths, code
% ---------------------------------------------------------------------
\usepackage{graphicx}
\graphicspath{{figures/}}
\usepackage{booktabs, tabularx, longtable, multirow, array}
\usepackage{enumitem}
\usepackage[font=small, labelfont=bf]{caption}
\usepackage{subcaption}
\usepackage{amsmath, amssymb}
\usepackage{siunitx}
\sisetup{group-separator={,}, group-minimum-digits=4, per-mode=symbol,
  mode=match}                     % numbers and units in the text font
\usepackage{listings}
\usepackage{pdfpages}
\usepackage[en-GB]{datetime2}

\usepackage{tikz}
\usetikzlibrary{positioning, arrows.meta, calc, fit, backgrounds,
  shapes.geometric, shapes.misc}
\usepackage{pgfplots}
\pgfplotsset{compat=1.18}
% plot colours: brand palette, high contrast (brand guidelines p. 37)
\pgfplotscreateplotcyclelist{solent}{
  {SolentRed, mark=*}, {Navy, mark=square*}, {Teal, mark=triangle*},
  {Orange, mark=diamond*}, {SkyBlue, mark=pentagon*}}
\pgfplotsset{every axis/.append style={cycle list name=solent}}
\usepackage{pgfplotstable}
\usepackage{pgfgantt}
\usepackage{tcolorbox}
\tcbuselibrary{skins, breakable}

% ---------------------------------------------------------------------
% Headers and footers: page number bottom centre (Handbook, App. A)
% ---------------------------------------------------------------------
\iffyp@poster\else
\usepackage[headsepline=false]{scrlayer-scrpage}
\clearpairofpagestyles
\cfoot*{\pagemark}
\pagestyle{scrheadings}
\fi

% ---------------------------------------------------------------------
% Draft-mode helpers: todo notes, watermark, header stamp
% ---------------------------------------------------------------------
\iffinalmode
  \usepackage[disable]{todonotes}
\else
  \usepackage[textsize=scriptsize, color=SkyBlue!30, bordercolor=SkyBlue]{todonotes}
  % todo notes are not counted as words
  \NewCommandCopy\fyp@todo\todo
  \renewcommand\todo[2][]{\directlua{fyp_wc.paused = fyp_wc.paused + 1}%
    \fyp@todo[#1]{#2}\directlua{fyp_wc.paused = fyp_wc.paused - 1}}
  \usepackage{draftwatermark}
  \SetWatermarkText{DRAFT}
  \SetWatermarkColor[gray]{0.9}
  \SetWatermarkScale{1.2}
  \iffyp@poster\else\ohead*{\scriptsize\textcolor{SolentRed}{\textbf{DRAFT}}\ compiled
    \DTMnow\ \textbar\ main body: \fyp@lastwords\ words
    (\iffyp@maxmode maximum\else target\fi\ \fyp@target)}\fi
\fi

% ---------------------------------------------------------------------
% References: Harvard Solent (Solent University's interpretation of
% BS ISO 690), see https://libguides.solent.ac.uk/HarvardSolent
% Same configuration as the SEF639 Project Handbook.
% ---------------------------------------------------------------------
\usepackage[
  backend=biber,
  style=ext-authoryear,
  sorting=nyt,
  maxbibnames=3, minbibnames=1,   % four or more authors -> et al.
  maxcitenames=3, mincitenames=1,
  giveninits=true,
  uniquename=false, uniquelist=false,
  dashed=false,
  date=year,
  urldate=long,
  isbn=false, doi=false, eprint=false
]{biblatex}

% In-text: (Johns 2014) / Johns (2014); several sources separated by ';'
\DeclareDelimFormat{nameyeardelim}{\addspace}
\DeclareDelimFormat{finalnamedelim}{\addspace\bibstring{and}\space}
\DeclareDelimFormat[bib]{finalnamedelim}{\addspace\bibstring{and}\space}
\DeclareDelimFormat{multicitedelim}{\addsemicolon\space}
% Reference list: JOHNS, R., 2014. \emph{Title}. 6th ed. London: Publisher
\DeclareDelimFormat[bib]{nameyeardelim}{\addcomma\space}
\DeclareFieldFormat{biblabeldate}{#1}
\DeclareNameAlias{sortname}{family-given/given-family}
\renewrobustcmd*{\bibinitdelim}{}
\AtBeginBibliography{\renewcommand*{\mkbibnamefamily}[1]{\MakeUppercase{#1}}}
\renewcommand*{\finentrypunct}{}
\DeclareFieldFormat{editortype}{#1}
% ... [viewed 8 October 2026]. Available from: URL
\newcommand*{\solentmonth}[1]{\ifcase#1\or January\or February\or March\or
  April\or May\or June\or July\or August\or September\or October\or
  November\or December\fi}
\DeclareFieldFormat{url}{Available from: \url{#1}}
\renewbibmacro*{url+urldate}{%
  \iffieldundef{url}{}{%
    \iffieldundef{urlyear}{}{%
      \setunit{\addspace}%
      \printtext[brackets]{viewed\space\thefield{urlday}\space
        \solentmonth{\thefield{urlmonth}}\space\thefield{urlyear}}}%
    \setunit{\adddot\space}\printfield{url}}}
% Split the list into 'References' (cited) and 'Bibliography' (consulted
% but not cited -- add those with \nocite{key} in main.tex)
\DeclareBibliographyCategory{cited}
\AtEveryCitekey{\addtocategory{cited}{\thefield{entrykey}}}

% ---------------------------------------------------------------------
% Cross-references and links
% ---------------------------------------------------------------------
\usepackage{xurl}
\urlstyle{same}
% Work-around for a varioref bug (v1.6i, 2026): define the language hook
% before varioref loads. See https://github.com/latex3/latex2e/issues/2112
\providecommand\extrasenglish{}
\usepackage{varioref}
\usepackage{hyperref}
\hypersetup{
  colorlinks=true,
  linkcolor=\iffinalmode SeaBlue\else SolentRed\fi,
  citecolor=SeaBlue,
  urlcolor=SeaBlue,
  breaklinks=true,
  pdfdisplaydoctitle=true,
  bookmarksnumbered=true
}
\usepackage[nameinlink]{cleveref}
\crefname{chapter}{Chapter}{Chapters}
\crefname{section}{Section}{Sections}
\crefname{subsection}{Section}{Sections}
\crefname{appendix}{Appendix}{Appendices}

% ---------------------------------------------------------------------
% Notation, abbreviations and glossary (entries in glossary-entries.tex)
% Sorted by LaTeX itself (\makenoidxglossaries): no external tool such as
% makeglossaries (Perl) or makeindex is needed.
% ---------------------------------------------------------------------
\usepackage[acronym, symbols, nonumberlist, nopostdot, nogroupskip]{glossaries}
\makenoidxglossaries

% ---------------------------------------------------------------------
% Headings required by the Handbook (Appendix A)
% ---------------------------------------------------------------------
\gappto\captionsenglish{%
  \renewcommand\contentsname{CONTENTS}%
  \renewcommand\listfigurename{FIGURES}%
  \renewcommand\listtablename{TABLES}%
  \renewcommand\bibname{References}%
  \renewcommand\refname{References}}
\setkomafont{disposition}{\color{Navy}\sffamily\bfseries}

% =====================================================================
% Your details (set in main.tex)
% =====================================================================
\newcommand\studentname[1]{\gdef\fyp@student{#1}}
\newcommand\studentid[1]{\gdef\fyp@studentid{#1}}
\newcommand\award[1]{\gdef\fyp@award{#1}}
\newcommand\course[1]{\gdef\fyp@course{#1}}
\newcommand\projecttitle[1]{\gdef\fyp@title{#1}}
\newcommand\supervisor[1]{\gdef\fyp@supervisor{#1}}
\newcommand\academicyear[1]{\gdef\fyp@year{#1}}
\newcommand\reportdate[1]{\gdef\fyp@date{#1}}
% \wordtarget{N}: aim for N words, +/-10% (Final Report)
% \wordlimit{N}:  N words is the maximum (Progress Report)
% \wordguide{N}:  N words is a recommendation only (Poster)
\newcommand\wordtarget[1]{\gdef\fyp@target{#1}\global\fyp@maxmodefalse
  \global\fyp@guidemodefalse}
\newcommand\wordlimit[1]{\gdef\fyp@target{#1}\global\fyp@maxmodetrue
  \global\fyp@guidemodefalse}
\newcommand\wordguide[1]{\gdef\fyp@target{#1}\global\fyp@maxmodetrue
  \global\fyp@guidemodetrue}
\newif\ifwordcountontitle
\studentname{Student Name}\studentid{00000000}\award{BEng (Hons)}
\course{Course Name}\projecttitle{Project Title}\supervisor{Supervisor}
% academic year and date follow the date of compilation unless set in
% main.tex: from September the new academic year (e.g. 2026--27)
\ifnum\month>8
  \edef\fyp@year{\the\year--\expandafter\@gobbletwo\the\numexpr\year+1\relax}
\else
  \edef\fyp@year{\the\numexpr\year-1\relax--\expandafter\@gobbletwo\the\year}
\fi
\reportdate{\solentmonth{\month}\ \the\year}
\wordtarget{10000}

\AtBeginDocument{\hypersetup{pdftitle={\fyp@title},
  pdfauthor={\fyp@student}, pdfsubject={SEF639 \fyp@doctype}}}

% =====================================================================
% Main body and word count
% =====================================================================
\directlua{dofile(kpse.find_file("setup/wordcount.lua") or "setup/wordcount.lua")
  fyp_wc.unit = "c@\fyp@unit"}
\providecommand\fyp@lastwords{--}   % value from the previous run (.aux)

% Headers and footers are built in the output routine: do not count them
\newcommand\fyp@pause{\directlua{fyp_wc.paused = fyp_wc.paused + 1}}
\newcommand\fyp@resume{\directlua{fyp_wc.paused = fyp_wc.paused - 1}}
\AtBeginDocument{\output=\expandafter{\expandafter\fyp@pause\the\output\fyp@resume}}

% Remember chapter/section titles for the word-count report
\expandafter\pretocmd\csname add\fyp@unit tocentry\endcsname%
  {\protected@edef\fyp@tmp{#2}\expandafter\global\expandafter\let\csname fyp@chaptitle@#1\endcsname\fyp@tmp}{}{}

\newcommand\startmainmatter{%
  \clearpage\pagenumbering{arabic}% main body starts on page 1
  \par\directlua{fyp_wc.active = true}}

\newcommand\stopmainmatter{%
  \par\directlua{fyp_wc.active = false}%
  \edef\fyp@words{\directlua{tex.sprint(fyp_wc.total)}}%
  \immediate\write\@mainaux{\gdef\string\fyp@lastwords{\fyp@words}}%
  \typeout{^^J*** Main body word count: \fyp@words\space
    (\iffyp@maxmode maximum\else target\fi\space\fyp@target) ***^^J}%
  \iffyp@guidemode
    \ifnum\fyp@words>\fyp@target\relax
      \PackageWarningNoLine{SEF639}{The text has \fyp@words\space words,
        more than the recommended \fyp@target. Consider replacing text
        with figures}\fi
  \else\iffyp@maxmode
    \ifnum\fyp@words>\fyp@target\relax
      \PackageWarningNoLine{SEF639}{Main body has \fyp@words\space words,
        more than the maximum of \fyp@target. Text beyond the maximum
        will be disregarded}\fi
  \else
    \ifnum\numexpr10*\fyp@words<\numexpr9*\fyp@target\relax
      \PackageWarningNoLine{SEF639}{Main body has \fyp@words\space words,
        more than 10\% below the target of \fyp@target}\fi
    \ifnum\numexpr10*\fyp@words>\numexpr11*\fyp@target\relax
      \PackageWarningNoLine{SEF639}{Main body has \fyp@words\space words,
        more than 10\% above the target of \fyp@target}\fi
  \fi\fi}

% Word-count report (draft mode only, printed after the appendices)
\newcommand\FYPwcrow[2]{%
  #1 & \ifcsname fyp@chaptitle@#1\endcsname
         \csname fyp@chaptitle@#1\endcsname\fi & \num{#2}\\}
\newcommand\wordcountreport{%
  \iffinalmode\else
  \clearpage
  \fyp@addunnumbered{Word Count Report (draft only)}
  This page is only printed in draft mode. It counts the words of the main
  body as typeset: body text, headings and list items. Captions, tables,
  footnotes, equations, guidance boxes, front matter, references and
  appendices are not counted.
  \iffyp@maxmode The Handbook sets a maximum length; any text beyond the
  maximum is disregarded and not taken into account in the mark.
  \else The Handbook expects the main body to be within $\pm$10\% of the
  target.\fi
  \begin{center}
  \begin{tabular}{rlr}
    \toprule
    & \textbf{\fyp@unitname} & \textbf{Words}\\ \midrule
    \directlua{fyp_wc.report_rows()}
    \midrule
    & \textbf{Main body total} & \textbf{\num{\fyp@words}}\\
    \iffyp@maxmode
    & Maximum & \num{\fyp@target}\\
    \else
    & Target ($\pm$10\%) & \num{\fyp@target}
      (\num{\the\numexpr9*\fyp@target/10\relax}--\num{\the\numexpr11*\fyp@target/10\relax})\\
    \fi
    \bottomrule
  \end{tabular}
  \end{center}
  \fi}

% =====================================================================
% Title page -- layout from Appendix E of the SEF639 Project Handbook,
% filled in from the details in main.tex.
% Logo: not part of the template. Download the Solent University logo from
% the SEF639 module page on SOL and save it as figures/solent-logo.pdf (or
% .png). Until then a placeholder is shown.
% =====================================================================
\newcommand\maketitlepage{%
\begin{titlepage}
  \centering
  \IfFileExists{figures/solent-logo.pdf}{\includegraphics[width=0.3\linewidth]{solent-logo}}{%
  \IfFileExists{figures/solent-logo.png}{\includegraphics[width=0.3\linewidth]{solent-logo}}{%
    {\Large\sffamily\bfseries Southampton Solent University\par}
    \placeholder[0.6\linewidth]{Logo: download it from the SEF639 module page
      on SOL and save it as \texttt{figures/solent-logo.pdf}}}}
  \vspace*{0.08\textheight}

  {\Large School of Technology and Maritime Industries\par}
  \vspace{0.04\textheight}
  {\color{SolentRed}\Huge\sffamily\bfseries \fyp@award\ \fyp@course\par}
  \vspace{0.08\textheight}

  {\Large
  \begin{tabular}{rl}
    Academic Year: & \fyp@year\\[2pt]
    Student:       & \fyp@student\\[2pt]
    Student ID:    & \fyp@studentid\\
  \end{tabular}\par}
  \vspace{0.05\textheight}
  {\LARGE\sffamily\bfseries\color{Navy} \fyp@title\par}
  \vfill

  {\Large Supervisor: \fyp@supervisor\par}
  \vspace{4pt}
  {\Large Date: \fyp@date\par}
  \vfill

  {\large SEF639 Final Year Project -- \fyp@doctype\par}
  \ifwordcountontitle
    {\normalsize Word count (main body): \fyp@lastwords\par}
  \else\iffinalmode\else
    \vspace{6pt}
    {\small\color{SolentRed} DRAFT \textbar\ main body: \fyp@lastwords\ words
      (\iffyp@maxmode maximum \fyp@target\else target \fyp@target\ $\pm$10\%\fi)\par}
  \fi\fi
\end{titlepage}}

% =====================================================================
% Front and back matter helpers
% =====================================================================
\newcommand\printnotation{%
  \printnoidxglossary[type=symbols, sort=standard, title=NOTATION, toctitle=Notation]%
  \printnoidxglossary[type=\acronymtype, sort=letter, title=ABBREVIATIONS, toctitle=Abbreviations]%
  \printnoidxglossary[sort=letter, title=GLOSSARY, toctitle=Glossary]}

\newcommand\printreferences{%
  \printbibliography[heading=bibintoc, title=References, category=cited]%
  \printbibliography[heading=bibintoc, title=Bibliography, notcategory=cited]}

% =====================================================================
% Building blocks (see the example chapters for how to use them)
% =====================================================================

% Guidance boxes: shown in draft mode, removed completely in final mode.
% Delete them as you write your own text.
\iffinalmode
  \NewDocumentEnvironment{guidance}{O{Guidance} +b}{}{}
\else
  \NewDocumentEnvironment{guidance}{O{Guidance}}
    {\par\directlua{fyp_wc.paused = fyp_wc.paused + 1}%
     \begin{tcolorbox}[enhanced, breakable, colback=white,
      colframe=SkyBlue, colbacktitle=SkyBlue!40, coltitle=Navy, fonttitle=\sffamily\bfseries\small,
      fontupper=\iffyp@poster\scriptsize\else\small\fi, title={#1},
      boxrule=0.6pt, arc=1mm,
      before upper={\setlength{\parskip}{4pt}}]}
    {\end{tcolorbox}\par\directlua{fyp_wc.paused = fyp_wc.paused - 1}}
\fi

% Grey placeholder box, e.g. for a missing figure or logo
% (in final mode a placeholder that is still there causes a warning)
\newcommand\placeholder[2][\linewidth]{%
  \iffinalmode\PackageWarning{SEF639}{A placeholder box is still in the
    final version. Replace it with your own content}\fi
  \begin{tcolorbox}[width=#1, colback=black!10, colframe=black!40,
    boxrule=0.5pt, arc=0pt, halign=center, valign=center,
    fontupper=\small\sffamily\color{black!60}, before skip=4pt, after skip=4pt]
  #2\end{tcolorbox}}

% Evidence PDF in an appendix, e.g.
%   \evidence[section]{GRA01 risk assessment}{evidence/gra01.pdf}{note}
% If the file exists: heading (#1 = section/subsection/chapter, default
% section) and all pages of the PDF. If it is missing: in draft mode the
% heading and a placeholder with the note; in final mode nothing at all,
% only a warning in the log.
\NewDocumentCommand\evidence{O{section} m m m O{0.85}}{%
  \IfFileExists{#3}{%
    \csname #1\endcsname{#2}%
    \includepdf[pages=-, scale=#5, pagecommand={}]{#3}}{%
    \iffinalmode
      \PackageWarningNoLine{SEF639}{Evidence file #3 not found:
        '#2' is left out of the final version}%
    \else
      \csname #1\endcsname{#2}%
      \placeholder{#4\\[2pt] save it as \texttt{#3}}%
    \fi}}

% Objectives O1, O2 ... that can be referenced with \cref
\newlist{objectives}{enumerate}{1}
\setlist[objectives]{label=\textbf{O\arabic*}, ref=O\arabic*, leftmargin=*}
\crefname{objectivesi}{objective}{objectives}
\Crefname{objectivesi}{Objective}{Objectives}

% Requirements R1, R2 ... : use \req{req:label} in the first column of a
% (non-tabularx) table; refer to them later with \cref{req:label}
\newcounter{requirement}
\renewcommand\therequirement{R\arabic{requirement}}
\newcommand\req[1]{\refstepcounter{requirement}\label{#1}\textbf{\therequirement}}
\crefname{requirement}{requirement}{requirements}
\Crefname{requirement}{Requirement}{Requirements}

% Test outcomes
\newcommand\pass{\textcolor{Teal}{\textbf{Pass}}}
\newcommand\fail{\textcolor{SolentRed}{\textbf{Fail}}}
\newcommand\partpass{\textcolor{Orange}{\textbf{Partial}}}

% Risk register: \riskrating{likelihood 1-5}{severity 1-5}
% prints the product and colours the table cell
\newcommand\riskrating[2]{%
  \edef\fyp@risk{\the\numexpr#1*#2\relax}%
  \ifnum\fyp@risk<5 \cellcolor{RiskLow}\fyp@risk\ (Low)%
  \else\ifnum\fyp@risk<10 \cellcolor{RiskMedium}\fyp@risk\ (Medium)%
  \else\ifnum\fyp@risk<16 \cellcolor{RiskHigh}\fyp@risk\ (High)%
  \else\cellcolor{RiskVeryHigh}\fyp@risk\ (Very high)\fi\fi\fi}

% Weighted decision matrix (use a plain tabular, not tabularx):
%   \resetscores{A,B,C}       at the start of the table
%   \weight{0.3}              prints the weight of the current criterion
%   \score{A}{4}              prints the score and adds weight x score to A
%   \weightsum, \scoretotal{A} in the last row
\newcommand\resetscores[1]{%
  \gdef\fyp@wsum{0}\gdef\fyp@w{0}%
  \@for\fyp@opt:=#1\do{\expandafter\gdef\csname fyp@tot@\fyp@opt\endcsname{0}}}
\newcommand\weight[1]{\gdef\fyp@w{#1}\xdef\fyp@wsum{\fpeval{\fyp@wsum+#1}}#1}
\newcommand\score[2]{#2\expandafter\xdef\csname fyp@tot@#1\endcsname
  {\fpeval{\csname fyp@tot@#1\endcsname+\fyp@w*(#2)}}}
\newcommand\scoretotal[1]{\fpeval{round(\csname fyp@tot@#1\endcsname,2)}}
\newcommand\weightsum{\fpeval{round(\fyp@wsum,2)}%
  \ifdim\fpeval{abs(\fyp@wsum-1)}pt>0.001pt
    \ \textcolor{SolentRed}{(should be 1)}\fi}

% Gantt charts (pgfgantt): \begin{ganttchart}[fyp]{1}{30} ... \end{ganttchart}
%   \fypsemesters            the two title rows: Semester 1 (12 weeks),
%                            break (3), Semester 2 (15 weeks) = 30 columns
%   \ganttgroup{..}{a}{b}    a project stage (Proposal, Progress Report,
%                            Final Report, Engineering Project Day)
%   \ganttbar{..}{a}{b}      planned task (grey), with WBS number and hours
%   \ganttbar[fyp actual]{}{a}{b}  actual duration (blue), same row
%   \ganttmilestone{..}{a}   submission (red)
\tikzset{
  /pgfgantt/fyp/.style={
    x unit=0.37cm, y unit chart=0.5cm, y unit title=0.5cm,
    vgrid={*1{black!10, line width=0.3pt}}, hgrid={*1{black!10}},
    title label font=\scriptsize, bar label font=\scriptsize,
    group label font=\scriptsize\bfseries,
    title/.style={fill=black!10, draw=black!40},
    bar/.style={fill=black!40, draw=none}, bar height=0.32,
    bar top shift=0.1,
    group/.style={fill=Navy, draw=none}, group height=0.2,
    group top shift=0.4, group left shift=0, group right shift=0,
    group peaks height=0.12, group peaks tip position=0,
    group peaks width=0.15,
    milestone/.style={fill=SolentRed, draw=none},
    milestone label font=\scriptsize\itshape,
    today rule/.style={draw=SolentRed, thick, dashed},
    today label font=\scriptsize\color{SolentRed}},
  /pgfgantt/fyp actual/.style={bar/.style={fill=SkyBlue, draw=none},
    bar top shift=0.5},
}
\newcommand\fypsemesters{%
  \gantttitle{Semester 1 (weeks)}{12}\gantttitle{Break}{3}%
  \gantttitle{Semester 2 (weeks)}{15}\\
  \gantttitlelist{1,...,12}{1}\gantttitle{}{3}\gantttitlelist{1,...,15}{1}\\}

% TikZ styles for flowcharts (methodology, system diagrams)
\tikzset{
  fyp node/.style={draw=Navy, thick, text=Navy, align=center, font=\small,
    execute at begin node={\hyphenpenalty=10000},
    minimum height=2.4em, inner sep=4pt},
  fyp process/.style={fyp node, rectangle, fill=SkyBlue!40,
    text width=9em},
  fyp terminal/.style={fyp node, rounded rectangle, fill=black!10,
    text width=7em},
  fyp decision/.style={fyp node, diamond, aspect=2.2, fill=Yellow!40,
    text width=6.5em, inner sep=1pt},
  fyp io/.style={fyp node, trapezium, trapezium left angle=72,
    trapezium right angle=108, trapezium stretches body,
    fill=black!10, text width=8em},
  fyp data/.style={fyp node, cylinder, shape border rotate=90,
    aspect=0.15, fill=black!10, text width=6em},
  fyp deliverable/.style={fyp node, rectangle, double, fill=SkyBlue!60,
    text width=9em},
  fyp arrow/.style={-{Stealth[length=2.5mm]}, thick, Navy},
  fyp group/.style={draw=black!40, dashed, rounded corners, inner sep=8pt}
}

% Code listings
\lstdefinestyle{fyp}{
  basicstyle=\ttfamily\footnotesize, numbers=left,
  numberstyle=\tiny\color{black!50}, frame=single, rulecolor=\color{black!25},
  keywordstyle=\color{SeaBlue}\bfseries, commentstyle=\color{Teal},
  stringstyle=\color{SolentRed}, breaklines=true, showstringspaces=false,
  captionpos=b, tabsize=2}
\lstset{style=fyp}

\makeatother

\wordguide{800}                  % recommendation, not a Handbook limit

\makeatletter
% ---- page: A2, landscape or portrait -------------------------------
\iffyp@portrait
  \usepackage[paperwidth=420mm, paperheight=594mm, margin=15mm]{geometry}
\else
  \usepackage[paperwidth=594mm, paperheight=420mm, margin=15mm]{geometry}
\fi
\pagestyle{empty}
\setlength{\parskip}{0.4em}
\setlength{\parindent}{0pt}
\color{Navy}
\usepackage{qrcode}

% the poster counts all text in its panels (they are boxes)
\directlua{fyp_wc.groups = { [""] = true, vbox = true, vtop = true }}

% ---- details only the poster needs ---------------------------------
\newcommand\fyp@paneltitle[1]{\csname fyp@pt@#1\endcsname}
\newcommand\paneltitle[2]{\expandafter\gdef\csname fyp@pt@#1\endcsname{#2}}
\paneltitle{intro}{Introduction and Aims}
\paneltitle{method}{Methodology}
\paneltitle{design}{Design and Implementation}
\paneltitle{results}{Results}
\paneltitle{conclusions}{Conclusions}
\paneltitle{references}{References}
\newcommand\acknowledgements[1]{\gdef\fyp@ack{#1}}
\newcommand\aistatement[1]{\gdef\fyp@ai{#1}}
\newcommand\qrlink[2]{\gdef\fyp@qrurl{#1}\gdef\fyp@qrtext{#2}}
\acknowledgements{}\aistatement{No AI tools were used in the preparation of this poster.}
\let\fyp@qrurl\@empty

% ---- styles ----------------------------------------------------------
\tcbset{
  fyp panel/.style={enhanced, colback=white, colframe=SolentRed,
    colbacktitle=SolentRed, coltitle=white, boxrule=1.2pt, arc=0pt,
    fonttitle=\sffamily\bfseries\LARGE, toptitle=3mm, bottomtitle=3mm,
    left=6mm, right=6mm, top=5mm, bottom=5mm, fontupper=\color{Navy},
    before upper={\setlength{\parskip}{0.4em}}},
  fyp plain/.style={enhanced, frame hidden, colback=white, boxrule=0pt,
    left=0pt, right=0pt, top=0pt, bottom=0pt, fontupper=\color{Navy}},
}

% ---- header ------------------------------------------------------------
\newcommand\fyp@logo{%
  \IfFileExists{figures/solent-logo.pdf}{\includegraphics[width=60mm]{solent-logo}}{%
  \IfFileExists{figures/solent-logo.png}{\includegraphics[width=60mm]{solent-logo}}{%
    \begin{tcolorbox}[width=60mm, height=60mm, colback=black!10,
      colframe=black!40, boxrule=0.5pt, arc=0pt, halign=center,
      valign=center, fontupper=\scriptsize\sffamily\color{black!60}]
      Logo: download it from the SEF639 module page on SOL and save it
      as figures/solent-logo.pdf
    \end{tcolorbox}}}}
\newcommand\fyp@header{%
  \begin{tcolorbox}[fyp plain, sidebyside, lefthand width=70mm, sidebyside gap=12mm]
    \fyp@logo
  \tcblower
    {\fontsize{54}{62}\selectfont\sffamily\bfseries\color{Navy}\fyp@title\par}
    \vspace{4mm}
    {\Large\sffamily\color{SolentRed}\bfseries\fyp@student\par}
    {\large \fyp@award\ \fyp@course\ \textbar\ Supervisor: \fyp@supervisor\par}
    {\large School of Technology and Maritime Industries \textbar\ SEF639
      Final Year Project \fyp@year\par}
  \end{tcolorbox}%
  \par\vspace{2mm}{\color{SolentRed}\rule{\linewidth}{2pt}}}

% ---- footer: acknowledgements, AI statement, QR code ------------------
\newcommand\fyp@footer{%
  \begin{tcolorbox}[fyp plain, fontupper=\normalsize\color{Navy}]
  \ifx\fyp@qrurl\@empty\else
    \begin{minipage}[c]{35mm}\centering
      \qrcode[height=30mm]{\fyp@qrurl}
    \end{minipage}\hspace{4mm}%
    \begin{minipage}[c]{60mm}\raggedright\small\fyp@qrtext\end{minipage}%
    \hspace{8mm}%
  \fi
  \begin{minipage}[c]{\dimexpr\linewidth-110mm\relax}
    \ifx\fyp@ack\@empty\else\textbf{Acknowledgements:} \fyp@ack\par\fi
    \textbf{Use of AI:} \fyp@ai\par
    \iffinalmode\else
      {\small\color{SolentRed} DRAFT \textbar\ compiled \DTMnow\ \textbar\
        poster text: \fyp@lastwords\ words (recommended up to \fyp@target)\par}
    \fi
  \end{minipage}
  \end{tcolorbox}}

% ---- panel content ---------------------------------------------------
\newcommand\fyp@panel[1]{%
  \directlua{fyp_wc.active = true}\input{panels/#1}\par
  \directlua{fyp_wc.active = false}}
\newcommand\fyp@refpanel{%
  {\small\printbibliography[heading=none, category=cited]}}

% ---- the poster ------------------------------------------------------
% Header and footer span the full width; in between, the panels are
% stacked in columns (landscape: three, portrait: two). The last panel of
% each column stretches to the bottom. If a panel is too full, LaTeX
% reports an overfull box: shorten the text or make a figure smaller.
\newsavebox\fyp@hbox \newsavebox\fyp@fbox
\newlength\fyp@colh \newlength\fyp@colw
\newlength\fyp@gap \setlength\fyp@gap{10mm}
\newcommand\fyp@box[2][]{% #1 extra options, #2 panel key/file
  \begin{tcolorbox}[fyp panel, title=\fyp@paneltitle{#2}, #1]
    \ifcsname fyp@special@#2\endcsname\csname fyp@special@#2\endcsname
    \else\fyp@panel{\csname fyp@file@#2\endcsname}\fi
  \end{tcolorbox}}
\@namedef{fyp@file@intro}{introduction}
\@namedef{fyp@file@method}{methodology}
\@namedef{fyp@file@design}{design}
\@namedef{fyp@file@results}{results}
\@namedef{fyp@file@conclusions}{conclusions}
\@namedef{fyp@special@references}{\fyp@refpanel}
% One column: the panels in #1 (comma list) at their natural height, then
% panel #2, which stretches to the bottom of the column.
\newsavebox\fyp@colbox \newsavebox\fyp@testbox \newlength\fyp@rest
\newcommand\fyp@column[2]{%
  \sbox\fyp@colbox{\begin{minipage}[t]{\fyp@colw}\vspace{0pt}%
    \@for\fyp@key:=#1\do{\fyp@box[nobeforeafter]{\fyp@key}\par
      \nointerlineskip\vspace*{\fyp@gap}}%
  \end{minipage}}%
  \setlength\fyp@rest{\dimexpr\fyp@colh-\ht\fyp@colbox-\dp\fyp@colbox\relax}%
  % warn if the column is too full (measured without counting words)
  \directlua{fyp_wc.paused = fyp_wc.paused + 1}%
  \sbox\fyp@testbox{\begin{minipage}[t]{\fyp@colw}\vspace{0pt}%
    \fyp@box[nobeforeafter]{#2}\end{minipage}}%
  \directlua{fyp_wc.paused = fyp_wc.paused - 1}%
  \ifdim\dimexpr\ht\fyp@testbox+\dp\fyp@testbox\relax>\fyp@rest
    \PackageWarningNoLine{SEF639}{The column ending with panel '#2' is too
      full by \the\dimexpr\ht\fyp@testbox+\dp\fyp@testbox-\fyp@rest\relax.
      Shorten the text or make a figure smaller}%
  \fi
  \begin{minipage}[t][\fyp@colh][t]{\fyp@colw}\vspace{0pt}%
    \usebox\fyp@colbox\par\nointerlineskip
    \fyp@box[height=\fyp@rest, nobeforeafter]{#2}%
  \end{minipage}}
\newcommand\makeposter{%
  \sbox\fyp@hbox{\parbox[b]{\textwidth}{\fyp@header}}%
  \sbox\fyp@fbox{\parbox[b]{\textwidth}{\fyp@footer}}%
  \setlength\fyp@colh{\dimexpr\textheight-\ht\fyp@hbox-\dp\fyp@hbox
    -\ht\fyp@fbox-\dp\fyp@fbox-2\fyp@gap-2pt\relax}%
  \iffyp@portrait
    \setlength\fyp@colw{\dimexpr(\textwidth-\fyp@gap)/2\relax}%
  \else
    \setlength\fyp@colw{\dimexpr(\textwidth-2\fyp@gap)/3\relax}%
  \fi
  \begingroup\setlength\parskip{0pt}%
  \noindent\usebox\fyp@hbox\par\nointerlineskip\vspace*{\fyp@gap}%
  \noindent
  \iffyp@portrait
    \fyp@column{intro,method}{design}\hspace{\fyp@gap}%
    \fyp@column{results,conclusions}{references}%
  \else
    \fyp@column{intro}{method}\hspace{\fyp@gap}%
    \fyp@column{design}{results}\hspace{\fyp@gap}%
    \fyp@column{conclusions}{references}%
  \fi
  \par\nointerlineskip\vspace*{\fyp@gap}\noindent\usebox\fyp@fbox\par
  \endgroup
  \stopmainmatter}
\makeatother
